\section{Conclusion}\label{sec:conclusion}

We proved that the \(p^3\)-case of the fixed-exponent Jacobson problem admits finite equational certificates for every prime \(p\).  Brandenburg's period criterion covers \(p=2\), \(p=3\), and \(p\equiv 2 \pmod 3\).  For the remaining primes \(p\equiv 1 \pmod 3\), the cyclic norm argument proves the two monomial constructive Wedderburn statements required by Brandenburg's reduction.

The essential new step is the normalization of a Frobenius-skew unit \(b\).  On each irreducible cubic component, the element \(b^3\) is fixed by the cycle; a finite norm section writes it as \(u\sigma(u)\sigma^2(u)\), and then \(u^{-1}b\) has cube \(1\).  Since \(T^3-1\) splits in the remaining congruence class, reducedness makes this normalized element central and eliminates the cubic component.  The proof gives a concrete route from the semantic commutativity theorem to explicit polynomial certificates for each fixed exponent \(p^3\).
